% ============================================================
%  EXAMEN — TEMA 1 (PRIMER LAPSO) — MODELO B
%  Números Racionales (Q): Definiciones, Fracciones Equivalentes,
%  Propiedades de la Potenciación, Operaciones, Potenciación y Ecuaciones.
%  Curso: 2do Año — CELAM
%  Autor: Ing. Gabriel Astudillo
%  Basado en el examen real (examen_real_t1l1.md).
%  Nota: SOLO el enunciado. No incluye clave de respuestas.
% ============================================================

\documentclass[12pt, letterpaper]{article}

% ── Paquetes ─────────────────────────────────────────────────
\usepackage{mathwizards-guia}
\mwguia{Examen — Tema 1 (Modelo B)}{Ing. Gabriel Astudillo $\cdot$ 2do Año}

% ── Comandos propios de este examen ──────────────────────────
\newcommand{\Q}{\mathbb{Q}}

\begin{document}
% ============================================================

% ── Portada / Título ─────────────────────────────────────────
\mwportada{Examen — Tema 1 (Modelo B)}{Números Racionales ($\Q$)}{Definiciones $\cdot$ Fracciones Equivalentes $\cdot$ Potenciación $\cdot$ Operaciones $\cdot$ Ecuaciones}

\vspace{4pt}
\begin{cuadroinfo}
  \small
  \textbf{Nombre:} \rule{0.55\linewidth}{0.4pt} \quad \textbf{Fecha:} \rule{0.15\linewidth}{0.4pt} \\[6pt]
  \textbf{Parámetros a considerar en la evaluación:}
  \begin{itemize}[leftmargin=*]
    \item Duración de la evaluación: \textbf{90 minutos}.
    \item El estudiante debe tener su lápiz y borrador.
    \item Ponderación total: \textbf{20 puntos}.
    \item Prohibido el uso de teléfonos o calculadoras.
    \item El estudiante debe permanecer en su puesto durante la evaluación.
  \end{itemize}
\end{cuadroinfo}

\vspace{8pt}

% ============================================================
\section{Definiciones y Ejemplos (4 puntos)}
% ============================================================

\begin{enumerate}[leftmargin=*]

  \item Defina y dé un ejemplo de los siguientes términos. \hfill \textbf{(Valor: 2 c/u)}
        \begin{enumerate}[label=\alph*), leftmargin=*]
          \item Máximo común divisor (M.C.D.).
          \item Números racionales.
        \end{enumerate}

\end{enumerate}

% ============================================================
\section{Fracciones Equivalentes (2 puntos)}
% ============================================================

\begin{enumerate}[leftmargin=*, resume]

  \item Escriba 4 pares de fracciones equivalentes entre sí (dos por amplificación y dos por simplificación). \hfill \textbf{(Valor: 2 puntos)}

\end{enumerate}

% ============================================================
\section{Propiedades de la Potenciación (2 puntos)}
% ============================================================

\begin{enumerate}[leftmargin=*, resume]

  \item Mencione 4 propiedades de la potenciación y haga un ejemplo de cada una. \hfill \textbf{(Valor: 2 puntos)}

\end{enumerate}

% ============================================================
\section{Operaciones con Números Racionales (4 puntos)}
% ============================================================

\begin{enumerate}[leftmargin=*, resume]

  \item Resuelva las siguientes operaciones con números racionales. \hfill \textbf{(Valor: 1 c/u)}
        \begin{enumerate}[label=\alph*), leftmargin=*]
          \item $\dfrac{2}{5} \cdot \left(\dfrac{6}{5} - \dfrac{3}{5}\right) =$
          \item $-\dfrac{3}{2} \cdot \left(\dfrac{7}{4} + \dfrac{5}{3}\right) =$
          \item $\dfrac{3}{4} + \dfrac{4}{5} \cdot \left(\dfrac{1}{3} - \dfrac{3}{2}\right) =$
          \item $-\dfrac{7}{4} - \dfrac{3}{2} \cdot \left(\dfrac{2}{3} + \dfrac{5}{4}\right) =$
        \end{enumerate}

\end{enumerate}

\newpage

% ============================================================
\section{Potenciación (4 puntos)}
% ============================================================

\begin{enumerate}[leftmargin=*, resume]

  \item Resuelva aplicando las propiedades de la potencia. Exprese el resultado \emph{sin exponentes negativos}. \hfill \textbf{(Valor: 1 c/u)}
        \begin{enumerate}[label=\alph*), leftmargin=*]
          \item $\left(\dfrac{18\,a^6 b^5}{24\,a^2 b^8}\right)^{2} =$
          \item $\left(\dfrac{4\,a^3 b^6}{8\,a^6 b^2}\right)^{-3} =$
          \item $\left(3a^2 b^{-3}\right)^2 \cdot \left(2a^{-1} b\right)^{4} =$
          \item $\dfrac{\left(4a^{-2} b^3\right)^{2}}{\left(2a^3 b^{-1}\right)} =$
        \end{enumerate}

\end{enumerate}

% ============================================================
\section{Ecuaciones (4 puntos)}
% ============================================================

\begin{enumerate}[leftmargin=*, resume]

  \item Resuelva las siguientes ecuaciones. \hfill \textbf{(Valor: 2 c/u)}
        \begin{enumerate}[label=\alph*), leftmargin=*]
          \item $3 + \dfrac{4}{5}x + \dfrac{5}{2}x = 8$
          \item $\dfrac{1}{2}y - \dfrac{7}{3}y + 6 = \dfrac{4}{5}$
        \end{enumerate}

\end{enumerate}

\vspace{10pt}
\begin{cuadroinfo}
  \small
  \textbf{Recuerda:} respeta la jerarquía de las operaciones, simplifica toda fracción hasta su forma irreductible y, en las ecuaciones, sustituye el valor obtenido para verificar la igualdad. \textquestiondown Mucho éxito!
\end{cuadroinfo}

\end{document}
